% p015_a 的电脑重画（安培力有效长度：弧线AB、直线AB、闭合圆环、与电容器相连的圆环）
\begin{tikzpicture}[line width=0.7pt, >={Stealth[length=2.2mm]}]
  % 左上：弧线 A→B
  \draw (1.2,5) .. controls (2.3,6.5) and (4.2,6.9) .. (4.2,5);
  \node at (0.6,4.55) {$A$};
  \node[right] at (4.2,4.85) {$B$};
  \node at (2.5,4.05) {$L_{\text{有效}}=\overline{AB}$};
  % 右上：直线 A→B
  \draw (6.4,5.8) -- (8.4,4.8);
  \node at (6.05,5.5) {$A$};
  \node[right] at (8.4,4.8) {$B$};
  \node at (6.8,4.05) {$L_{\text{有}}=\overline{AB}$};
  % 左下：闭合圆环
  \draw (2.7,2.4) circle (0.62);
  \node at (2.6,1.25) {$L=0$};
  % 右下：与电容器相连的圆环（内有 d 和一条虚线）
  \coordinate (K) at (6.6,2.45);
  \coordinate (WL) at ($(K)+(200:0.72)$);
  \coordinate (WR) at ($(K)+(-20:0.72)$);
  \draw (K) circle (0.72);
  \draw[dashed] (WL) -- (WR);
  \node at ($(K)+(0,0.22)$) {$d$};
  \draw[->] ($(K)+(70:0.72)$) arc (70:88:0.72);
  \draw[->] ($(K)+(-40:0.72)$) arc (-40:-77:0.72);
  \draw (WL) -- ++(-0.8,0) |- (6.5,1.3);
  \draw (WR) -- ++(0.8,0) |- (6.7,1.3);
  \draw (6.5,1.0) -- (6.5,1.6);
  \draw (6.7,1.0) -- (6.7,1.6);
  \node[anchor=west] at (8.75,1.6) {$L_{\mbox{\unclear{上}}}=d$};
  \node[anchor=west] at (8.75,0.45) {$L_{\mbox{\unclear{下}}}=d$};
\end{tikzpicture}
